\section{四川的一张纸，河北的一道河}

993年王小波在成都平原起事，次年李顺继领并一度占据成都。宋史所载“均贫富”是理解这场动员的重要线索，却不能被扩大为全体参与者共享的完整纲领；盐茶、赋役、军费和地方控制的压力，需要同成都平原的社会结构分别考察。宋军的平定恢复了主要城市控制，但不能由此推断战后家庭、土地和交通已立刻复原。

1012--1013年朝廷向江淮等地推广早熟稻种，后世常以“占城稻”概括；1023年又在益州设交子务，收取民间交子、发行官交子。前者见《宋史》食货志和会要的农田、稻种条，后者见食货志与交子条，均可证明国家曾作出相应安排。品种来源、传播速度和产量不能从一条诏令推及全国；交子的发行额、准备金和实际流通范围也须区别。两件事共同指向的并非自动发生的“商业革命”，而是朝廷试图在区域粮食与信用网络中取得可管理的接口。

1047--1048年贝州王则起事，1048年黄河又在商胡决口、改道北流。前者的宗教组织、参与者构成和伤亡数字受材料限制，不能扩大为全国性社会运动；后者的受灾人口、河道持续时间和河工成效，亦须以古河道沉积、堤防遗址和地方材料检验。本纪与会要可告诉我们朝廷何时得知、何时议工赈，却不能替河北每一处聚落发言。对边防地区而言，河流改道同时改变了耕地、道路、驻军与运输的条件。

\section{国家进入市场与知识的通道}

1057年礼部进士试、1072年市易法、1084年司马光进呈《资治通鉴》、1092年苏颂等完成新仪象台，分别把考试、市场、史书和授时带入国家能力的不同侧面。欧阳修主持科举的榜单确实塑造了新的士人网络，却不能以“一榜定文风”遮蔽长期变化。市易务的机构可见于会要，市场价格与商人处境仍需另证。《资治通鉴》进书表和《长编》说明它进入政治世界的方式，不使其成为唯一的历史解释；新仪象台的技术文本可靠，实物早佚，现代复原也不应被倒写为当年每一次运行的证据。

1103年崇宁兴学、三舍法扩展，1104年方田均税等丈量与赋税整顿重新推进。学校设置和方田条文可由《宋史》选举、食货志及《宋会要辑稿》确认，但入学人数、实际取士和丈量面积不能由制度名目推出。学校、科举和地方士人的路径有地区差异；丈量则使州县更需要土地信息，也可能把申报、隐田与税负冲突压到胥吏、豪强和小户之间。

1120年方腊在两浙起事，次年被童贯等军平定。镇压方材料保留了起事扩展和军事结果，也最容易放大口号、教派和斩获数。较稳妥的认识是，徭役、征敛、乡里与宗教网络在东南相互交织，威胁到漕运与税源；平乱恢复了官府据点，却抽走了北方战略所需的精锐，并将战后负担留在两浙。城市、市场和乡村从未只是京师财政的被动背景，它们既供给王朝，也会在征发过重、信息失真或灾害来临时显出自己的限度。
